%!TEX root = ./main.tex

\section[Today]{How Communication Works Today}

% SECTION TIMING: 2:20--14:10 (5 frames): 130 + 190 + 130 + 90 + 170 = 710 s.
% Flex frame if running late: "Why this design won" compresses to 60 s (say the two
% ideas, skip the formula).

% Teaching note (130 s): redraw Shannon's Fig. 1 and name the five boxes (in the
% original figure the channel is an unlabeled box; the text names the five parts:
% source, transmitter, channel, receiver, destination). Shannon's problem: reproduce,
% at the destination, the message chosen at the source -- exactly or approximately.
% Then read the quote slowly; it is the hinge of the whole lecture. Shannon did not
% say meaning is unimportant, only that it is not needed to design the pipe.
% Information = which message was chosen out of all possible ones, counted in bits.
\begin{frame}{Shannon 1948: copy the message}
  \vspace{1.2em}
  \centering
  % SOURCE: Shannon, Bell Syst. Tech. J. 27 (1948), Fig. 1 and Introduction.
  \begin{tikzpicture}[scale=0.74,transform shape,
      blk/.style={netnode,minimum width=2.1cm,minimum height=1.0cm,font=\small\bfseries}]
    \node[blk] (src) at (0,0) {Information\\source};
    \node[blk] (tx)  at (3.05,0) {Transmitter};
    \node[chnode,minimum width=1.8cm,minimum height=1.0cm,font=\small\bfseries] (ch)
      at (6.1,0) {Channel};
    \node[blk] (rx)  at (9.15,0) {Receiver};
    \node[blk] (dst) at (12.2,0) {Destination};
    \node[corenode,minimum width=1.8cm,minimum height=0.75cm,font=\small\bfseries]
      (noise) at (6.1,-1.45) {Noise};
    \draw[flow] (src) -- (tx);
    \draw[flow] (tx) -- (ch);
    \draw[flow] (ch) -- (rx);
    \draw[flow] (rx) -- (dst);
    \draw[flow] (noise) -- (ch);
  \end{tikzpicture}

  \vspace{1.0em}
  \uncover<2->{{\large Goal: an \textbf{exact copy} at the other end.}\par}

  \vspace{1.0em}
  \uncover<3->{{\itshape ``These semantic aspects of communication are irrelevant
    to the engineering problem.''}\par
    \vspace{0.2em}
    {\small\textcolor{softgray}{--- Shannon (1948)}}\par}

  \vspace{0.9em}
  \takeawayat{4-}{Meaning was left out on purpose:\\
    a bit is a bit, whatever it encodes.}
  \source{https://doi.org/10.1002/j.1538-7305.1948.tb01338.x}{C.\ E.\ Shannon,
    ``A Mathematical Theory of Communication,'' Bell Syst.\ Tech.\ J., 1948}
\end{frame}

% Teaching note (190 s): the suitcase analogy carries the slide. Compress = vacuum-pack
% the clothes (source coding; for video, a codec such as H.264 removes what is
% predictable inside each frame and between consecutive frames, so every bit left
% matters). Protect = bubble wrap (channel coding: add structured redundancy so the
% receiver can fix errors; modern examples are LDPC and polar codes -- the Wi-Fi
% lecture did not cover them, so just say "error-correcting code"). Send = the truck
% (modulation: the QAM and OFDM of the Wi-Fi lecture). The truck never opens the
% suitcase: each block is designed and upgraded on its own, and the halves agree on
% one interface only -- bits. The receiver undoes each step in reverse order.
% Fourth click plants FEEDBACK, which "Where it struggles" builds on: the receiver
% keeps reporting how the delivery is going, and the sender adjusts -- phone-call
% analogy: "you're breaking up", so you slow down and repeat. In real-time video
% (WebRTC) the receiver reports packet arrival times (TWCC) and the sender's
% congestion control (GCC) sets the encoder's target bitrate -- hence the arrow back
% to Compress. Wi-Fi itself also adapts at the Send step (it switches to slower,
% sturdier transmission settings); no need to say this unless asked.
\begin{frame}{Today's pipeline: like shipping a suitcase}
  \vspace{1.6em}
  \centering
  \begin{tikzpicture}[scale=0.72,transform shape,
      blk/.style={netnode,minimum width=2.3cm,minimum height=1.0cm,font=\small\bfseries}]
    \node[endnode] (src) at (0,0) {Video};
    \node[blk] (c) at (2.8,0) {Compress\\[-2pt]{\scriptsize\mdseries e.g., H.264}};
    \node[blk] (p) at (5.75,0) {Protect\\[-2pt]{\scriptsize\mdseries error-correcting code}};
    \node[blk] (s) at (8.7,0) {Send\\[-2pt]{\scriptsize\mdseries radio signal}};
    \node[chnode,minimum width=1.3cm,minimum height=1.0cm] (air) at (11.2,0) {Air};
    \node[endnode,minimum width=1.9cm] (rx) at (13.5,0) {Receiver};
    \draw[flow] (src) -- (c);
    \draw[flow] (c) -- (p);
    \draw[flow] (p) -- (s);
    \draw[flow] (s) -- (air);
    \draw[flow] (air) -- (rx);
    \node[iface,above=3pt of c,visible on=<2->] {vacuum-pack};
    \node[iface,above=3pt of p,visible on=<2->] {bubble wrap};
    \node[iface,above=3pt of s,visible on=<2->] {truck};
    \draw[flow,coreorange,dashed,visible on=<4->] (rx.south) -- ++(0,-0.9) -| (c.south);
    \node[font=\small\bfseries,text=coreorange,visible on=<4->] at (8.2,-1.05)
      {feedback: ``too much, slow down!''};
  \end{tikzpicture}

  \vspace{1.0em}
  \uncover<3->{{\large The truck never opens the suitcase.}\par}

  \vspace{0.5em}
  \uncover<4->{{\small\textcolor{softgray}{The receiver reports back, and the sender
    adjusts the video quality --- like ``you're breaking up'' on a phone call.}}\par}

  \vspace{1.0em}
  \takeawayat{5-}{Every block has one job:\\
    \textbf{deliver these bits exactly.}}
\end{frame}

% Teaching note (130 s): reveal one box at a time and do the arithmetic out loud.
% Raw: 1920 x 1080 pixels x 24 bits x 30 frames/s = 1.49 Gb/s. H.264 at 8 Mb/s: that
% is YouTube's recommended upload bitrate for 1080p SDR at 24/25/30 frames/s (YouTube
% Help, "YouTube recommended upload encoding settings"; 12 Mb/s at 48/50/60 frames/s)
% -- a realistic compressed rate, about 190x smaller than raw. Protection: the 802.11a
% 54 Mb/s mode uses a rate-3/4 convolutional code, i.e., +1/3 extra bits -> about
% 10.7 Mb/s on air, roughly 20% of an old 54 Mb/s Wi-Fi link, nonstop. The orange box
% is a TEACHING EXAMPLE: if B only needs a pedestrian count, 1 byte per frame x 30
% frames/s = 240 b/s. Q&A mechanism (unicast Wi-Fi frames): the 32-bit frame check
% sequence (a CRC) fails, the receiver drops the frame and sends no ACK, the sender
% times out and retries up to its retry limit; broadcast frames are not ACKed, and
% inside an A-MPDU only the failed subframe is resent. Sources: IEEE Std 802.11;
% Ergen, "IEEE 802.11 Tutorial," UC Berkeley, 2002.
\begin{frame}{One video stream, by the numbers}
  \vspace{1.4em}
  \centering
  % SOURCE (8 Mb/s): https://support.google.com/youtube/answer/1722171
  \begin{tikzpicture}[scale=0.8,transform shape,
      box/.style={netnode,minimum width=2.9cm,minimum height=1.25cm,font=\small\bfseries}]
    \node[box] (raw) at (0,0) {Raw 1080p video\\[1pt]{\large 1.5 Gb/s}};
    \node[box,visible on=<2->] (h264) at (3.9,0) {H.264\\[1pt]{\large 8 Mb/s}};
    \node[box,visible on=<3->] (fec) at (7.8,0) {+ protection\\[1pt]{\large 10.7 Mb/s}};
    \draw[flow,visible on=<2->] (raw) -- (h264);
    \draw[flow,visible on=<3->] (h264) -- (fec);
    \node[lbl,below=3pt of raw] {30 frames per second};
    \node[lbl,below=3pt of h264,visible on=<2->] {YouTube's 1080p rate};
    \node[lbl,below=3pt of fec,visible on=<3->] {+1/3 extra bits};
    \node[corenode,minimum width=6.4cm,minimum height=0.95cm,font=\small\bfseries,
      visible on=<4->] (b) at (3.9,-2.2) {But B only needs ``2 people'': $\approx$240 b/s};
  \end{tikzpicture}

  \vspace{0.6em}
  \raggedright
  \uncover<5->{\qa{What does the receiver check at the end?}
     {Only whether every bit is right. One wrong bit, and it is sent again.
      Nobody asks what the video shows.}}

  \vspace{0.2em}
  \source{https://support.google.com/youtube/answer/1722171}{YouTube recommended
    upload encoding settings; IEEE Std 802.11; orange box = teaching example}
\end{frame}

% Teaching note (90 s): two reasons this design won. (1) Divide and conquer:
% Shannon's separation theorem says compressing and protecting can be designed
% separately with no loss -- for one sender and one receiver, and only in the limit of
% very long messages. Exact conditions (Cover & Thomas, Elements of Information
% Theory, 2nd ed., Wiley 2006, Thm. 7.13.1): point-to-point, discrete memoryless
% channel, source satisfying the AEP, error probability -> 0 only as block length ->
% infinity. Lossy form: distortion D is achievable if R(D) < C and impossible if
% R(D) > C. Consequence: H.264 does not care whether it rides on Wi-Fi or 5G -- bits are
% a universal interface. (2) A speed limit: capacity C = B log2(1 + SNR); e.g. 20 MHz
% at 20 dB SNR gives about 133 Mb/s. It is a ceiling on reliable bits per second, not
% a calculator for a real radio's rate.
\begin{frame}{Why this design won}
  \vspace{1.4em}
  \begin{columns}[T,onlytextwidth]
    \column{0.48\textwidth}
      \centering
      {\Large\textbf{\textcolor{ranblue}{Divide and conquer}}\par}
      \vspace{0.8em}
      Compress and protect\\separately: \textbf{nothing lost}\\
      {\small\textcolor{softgray}{(long messages, one link)}}
    \column{0.48\textwidth}
      \centering
      \uncover<2->{%
      {\Large\textbf{\textcolor{coreorange}{A speed limit}}\par}
      \vspace{0.8em}
      Every channel has a\\\textbf{maximum reliable rate}\\
      {\small\textcolor{softgray}{$C = B \log_2(1+\mathrm{SNR})$}}}
  \end{columns}

  \vspace{2.2em}
  \takeawayat{3-}{Bits became a universal interface:\\
    any app over any network.}
  \source{https://doi.org/10.1002/j.1538-7305.1948.tb01338.x}{Shannon (1948);
    Cover \& Thomas, Elements of Information Theory, 2nd ed., Thm.\ 7.13.1}
\end{frame}

% Teaching note (170 s): the pictures carry the slide -- one frame per click.
% Real video calls do not just freeze when the signal weakens: WebRTC-style feedback
% adapts. The receiver keeps reporting when each packet arrived (TWCC feedback; in
% libwebrtc the report interval defaults to 100 ms within 50-250 ms -- from an earlier
% reading of the source, recheck before quoting), and the sender's congestion control
% (GCC) estimates the available bandwidth and lowers the encoder's target bitrate, so
% the picture gets blockier rather than stopping (second frame). But the loop always
% reacts after the fact: when wireless capacity drops suddenly, packets queue up at
% the access point before the sender slows down, so frames arrive late -- a latency
% spike -- and the receiver's playout can briefly freeze (WebRTC's own statistics
% count such freezes as freezeCount). That gap between "the channel changed" and "the
% sender reacted" is exactly what Zhuge (Meng et al., "Achieving Consistent Low
% Latency for Wireless Real-Time Communications with the Shortest Control Loop," ACM
% SIGCOMM 2022, doi:10.1145/3544216.3544225) attacks by shortening the control loop
% at the AP. Background if asked: a link with ONE fixed setting and no feedback shows
% the pure all-or-nothing "cliff" (digital TV freezing; old analog TV just got snowy)
% -- learned-JSCC papers often compare against that, but the fair baseline is a
% system with feedback. Theory behind "long messages, one link": at finite
% blocklength separation pays a penalty that joint coding can avoid (Kostina & Verdu,
% IEEE Trans. Inf. Theory 2013, doi:10.1109/TIT.2013.2238657); for a Gaussian source
% over an AWGN channel with matched bandwidth, uncoded transmission is optimal at
% every SNR at once, which no single digital rate can match (Gastpar, Rimoldi &
% Vetterli, IEEE Trans. Inf. Theory 2003, doi:10.1109/TIT.2003.810631).
\begin{frame}{Where it struggles}
  \vspace{0.8em}
  \centering
  \begin{tikzpicture}[scale=0.72,transform shape]
    \begin{scope}
      \signalbars{0.45}{2.3}{4}
      \node[anchor=west,font=\small\bfseries] at (1.65,2.45) {strong};
      \pic at (0,0) {streetframe};
      \node[font=\small\bfseries,text=softgreen,anchor=north] at (1.7,-0.1) {sharp};
    \end{scope}
    \begin{scope}[visible on=<2->]
      \signalbars{4.45}{2.3}{2}
      \node[anchor=west,font=\small\bfseries] at (5.65,2.45) {weaker};
      \pic at (4,0) {streetblocky};
      \node[font=\small\bfseries,text=ranblue,anchor=north,align=center] at (5.7,-0.1)
        {lower quality\\[-2pt]{\scriptsize\mdseries feedback cut the bitrate}};
    \end{scope}
    \begin{scope}[visible on=<3->]
      \signalbars{8.45}{2.3}{1}
      \node[anchor=west,font=\small\bfseries] at (9.65,2.45) {sudden drop};
      \pic at (8,0) {streetframe};
      \pic at (8,0) {late};
      \node[font=\small\bfseries,text=coreorange,anchor=north,align=center] at (9.7,-0.1)
        {late, may freeze\\[-2pt]{\scriptsize\mdseries feedback not there yet}};
    \end{scope}
  \end{tikzpicture}

  \vspace{0.5em}
  \raggedright
  \begin{itemize}[<+(3)->]\itemsep5pt
    \item \textbf{Feedback adapts --- one step late:} quality drops, delay spikes
    \item \textbf{Every bit counts the same:} sky = pedestrian
    \item \textbf{Receiver B only needs an answer}
  \end{itemize}

  \vspace{0.5em}
  \takeawayat{7-}{Today's video adapts, but after the fact ---\\
    and it always sends everything.}
  \source{https://doi.org/10.1145/3544216.3544225}{Feedback delay on wireless links:
    Z.\ Meng et al., ``Zhuge,'' ACM SIGCOMM 2022}
\end{frame}
